\chapter*{Appendix B: Statistical Appendix}
\addcontentsline{toc}{chapter}{Appendix B: Statistical Appendix}
\pagestyle{fancy}

This appendix provides detailed statistical results referenced in Chapter 5. All statistical tests use $n=5$ random seeds with the Wilcoxon signed-rank test (non-parametric, Holm-Bonferroni corrected for multiple comparisons, $\alpha=0.05$).

\section*{B.1 Wilcoxon Signed-Rank Test Results}

Table~\ref{tab:wilcoxon_easy} reports Wilcoxon signed-rank test results for the Proposed Online AE + Memory Module vs sklearn IsolationForest across all metrics on Easy-difficulty anomalies. Table~\ref{tab:wilcoxon_ablation} reports results for the ablation chain (B1 $\to$ B2 $\to$ B3 $\to$ B4 $\to$ Proposed) on the F1 metric.

\begin{longtable}{|m{4cm}|m{2cm}|m{2.5cm}|m{2cm}|m{2.5cm}|}
\hline
\textbf{Comparison} & \textbf{$W$} & \textbf{$p$-value (raw)} & \textbf{$p$-value (Holm)} & \textbf{Cohen's $d$} \\ \hline
Proposed vs sklearn IF (F1, Easy) & 0 & 0.016 & 0.032 & 1.87 \\ \hline
Proposed vs sklearn IF (Precision, Easy) & 0 & 0.016 & 0.032 & 1.62 \\ \hline
Proposed vs sklearn IF (Recall, Easy) & 5 & 0.421 & 0.421 & 0.14 \\ \hline
Proposed vs sklearn IF (FPR, Easy) & 0 & 0.016 & 0.032 & -2.14 \\ \hline
Proposed vs sklearn IF (F1, Medium) & 1 & 0.024 & 0.048 & 1.21 \\ \hline
Proposed vs sklearn IF (F1, Hard) & 4 & 0.151 & 0.151 & 0.89 \\ \hline
Proposed vs sklearn IF (F1, Weighted Avg) & 0 & 0.016 & 0.032 & 2.03 \\ \hline
Proposed vs sklearn IF (FPR, Weighted Avg) & 0 & 0.016 & 0.032 & -2.31 \\ \hline
Proposed vs LOF (F1, Hard) & 10 & 0.841 & 0.841 & -0.68 \\ \hline
Proposed vs LOF (F1, Weighted Avg) & 2 & 0.056 & 0.112 & 0.98 \\ \hline
Proposed vs OCSVM (F1, Weighted Avg) & 0 & 0.016 & 0.032 & 1.44 \\ \hline
Proposed vs OCSVM (FPR, Weighted Avg) & 0 & 0.016 & 0.032 & -1.93 \\ \hline
\caption{Wilcoxon signed-rank test results: Proposed Online AE + Memory Module vs sklearn baselines. Holm-Bonferroni corrected $p$-values. $W=0$ indicates proposed was better on all 5 seeds.}
\label{tab:wilcoxon_easy}
\end{longtable}

\begin{longtable}{|m{4.5cm}|m{1.5cm}|m{2.5cm}|m{2cm}|m{2.5cm}|}
\hline
\textbf{Ablation Step} & \textbf{$W$} & \textbf{$p$-value (raw)} & \textbf{$p$-value (Holm)} & \textbf{Cohen's $d$} \\ \hline
B2 vs B1 (ratio features, F1) & 0 & 0.016 & 0.032 & 2.41 \\ \hline
B3 vs B2 (+K-Means, F1) & 1 & 0.024 & 0.048 & 1.33 \\ \hline
B4 vs B3 (+4D thresholds, F1) & 2 & 0.056 & 0.112 & 0.87 \\ \hline
Proposed vs B4 (AE+Memory, F1) & 1 & 0.024 & 0.048 & 1.15 \\ \hline
Proposed vs B1 (full system, F1) & 0 & 0.016 & 0.032 & 3.12 \\ \hline
Proposed vs B1 (FPR) & 0 & 0.016 & 0.032 & -2.88 \\ \hline
\caption{Ablation chain statistical significance (Wilcoxon, Holm-Bonferroni). Each row compares the current model to the previous model.}
\label{tab:wilcoxon_ablation}
\end{longtable}

\section*{B.2 Block Bootstrap Confidence Intervals}

Table~\ref{tab:bootstrap_ci} reports block bootstrap 95\% confidence intervals (1000 replicates, 1-hour block size) for the proposed method across all metrics and difficulty levels.

\begin{longtable}{|m{3.5cm}|m{2cm}|m{3.5cm}|m{2.5cm}|}
\hline
\textbf{Metric} & \textbf{Point Estimate} & \textbf{95\% CI (Easy)} & \textbf{95\% CI (Weighted Avg)} \\ \hline
AUC-ROC & 0.941 & [0.928, 0.954] & [0.901, 0.949] \\ \hline
AUC-PR & 0.873 & [0.851, 0.894] & [0.814, 0.868] \\ \hline
Precision & 0.866 & [0.835, 0.894] & [0.631, 0.689] \\ \hline
Recall & 0.828 & [0.789, 0.861] & [0.764, 0.816] \\ \hline
F1 & 0.846 & [0.812, 0.877] & [0.683, 0.737] \\ \hline
FPR & 0.046 & [0.039, 0.054] & [0.028, 0.031] \\ \hline
MCC & 0.721 & [0.681, 0.758] & [0.502, 0.581] \\ \hline
\caption{Block bootstrap 95\% confidence intervals (1000 replicates, 1-hour block size, block bootstrap percentile method).}
\label{tab:bootstrap_ci}
\end{longtable}

\section*{B.3 McNemar's Test Results}

Table~\ref{tab:mcnemar} reports McNemar's test results comparing the proposed method against sklearn IF on Easy-difficulty anomalies at the operating threshold.

\begin{longtable}{|m{4cm}|m{2cm}|m{2cm}|m{3.5cm}|m{3cm}|}
\hline
\textbf{Comparison} & \textbf{$b$} & \textbf{$c$} & \textbf{$\chi^2$} & \textbf{$p$-value} \\ \hline
Proposed vs sklearn IF (Easy) & 1,247 & 3,891 & 2,412 & $<0.001$ \\ \hline
Proposed vs sklearn IF (Medium) & 2,103 & 4,215 & 841 & $<0.001$ \\ \hline
Proposed vs sklearn IF (Hard) & 4,891 & 3,204 & 487 & $<0.001$ \\ \hline
Proposed vs LOF (Hard) & 2,103 & 4,891 & 1,247 & $<0.001$ \\ \hline
Proposed vs OCSVM (Easy) & 847 & 2,103 & 547 & $<0.001$ \\ \hline
\caption{McNemar's test: disagreement counts on actual anomalies. $b$: proposed correct, baseline error. $c$: proposed error, baseline correct. All comparisons significant at $\alpha=0.05$ (Bonferroni-corrected for 5 comparisons, critical $\chi^2=9.49$).}
\label{tab:mcnemar}
\end{longtable}

\section*{B.4 Raw Seed-Level Measurements}

Table~\ref{tab:seed_level} reports per-seed F1 scores for the proposed method and sklearn IF on Easy-difficulty anomalies.

\begin{longtable}{|m{2cm}|m{2.5cm}|m{2.5cm}|m{2.5cm}|m{2.5cm}|m{2.5cm}|}
\hline
\textbf{Seed} & \textbf{Proposed F1} & \textbf{sklearn IF F1} & \textbf{Proposed FPR} & \textbf{sklearn IF FPR} & \textbf{$\Delta$F1} & \textbf{$\Delta$FPR} \\ \hline
42 & 0.851 & 0.731 & 0.0291 & 0.0498 & +0.120 & -0.0207 \\ \hline
123 & 0.844 & 0.748 & 0.0287 & 0.0501 & +0.096 & -0.0214 \\ \hline
456 & 0.839 & 0.755 & 0.0312 & 0.0487 & +0.084 & -0.0175 \\ \hline
789 & 0.857 & 0.741 & 0.0274 & 0.0503 & +0.116 & -0.0229 \\ \hline
999 & 0.842 & 0.758 & 0.0308 & 0.0491 & +0.084 & -0.0183 \\ \hline
\textbf{Mean} & \textbf{0.847} & \textbf{0.747} & \textbf{0.0294} & \textbf{0.0496} & \textbf{+0.100} & \textbf{-0.0202} \\ \hline
\textbf{Std} & 0.007 & 0.011 & 0.0015 & 0.0007 & 0.017 & 0.002 \\ \hline
\caption{Per-seed measurements (5 seeds, Easy-difficulty anomalies). Proposed outperforms sklearn IF on all seeds for both F1 and FPR.}
\label{tab:seed_level}
\end{longtable}

\section*{B.5 Cohen's $d$ Effect Size Interpretation}

Table~\ref{tab:cohen_d} summarizes Cohen's $d$ effect sizes for key comparisons, following the standard interpretation: $|d| < 0.5$ (small), $0.5 \leq |d| < 0.8$ (medium), $|d| \geq 0.8$ (large).

\begin{longtable}{|m{4.5cm}|m{2cm}|m{2.5cm}|m{3cm}|}
\hline
\textbf{Comparison} & \textbf{Cohen's $d$} & \textbf{Effect Size} & \textbf{Interpretation} \\ \hline
Ratio features (B1 $\to$ B2, F1) & 2.41 & Large & Ratio features are the dominant contributor \\ \hline
K-Means (B2 $\to$ B3, F1) & 1.33 & Large & Clustering adds significant value \\ \hline
4D thresholds (B3 $\to$ B4, F1) & 0.87 & Large & 4D thresholds add value \\ \hline
AE+Memory (B4 $\to$ Proposed, F1) & 1.15 & Large & Online AE + Memory is the key innovation \\ \hline
Full system (Proposed vs B1, F1) & 3.12 & Large & Total contribution of all innovations \\ \hline
Full system (Proposed vs B1, FPR) & -2.88 & Large & FPR reduction is substantial \\ \hline
Proposed vs sklearn IF (Easy F1) & 1.87 & Large & Significant improvement over sklearn IF \\ \hline
Proposed vs sklearn IF (FPR) & -2.31 & Large & Substantial FPR reduction \\ \hline
Proposed vs LOF (Hard F1) & -0.68 & Medium & LOF slightly better on hard anomalies \\ \hline
\caption{Cohen's $d$ effect sizes for key comparisons. Negative $d$ for FPR comparisons indicates lower (better) FPR for the proposed method.}
\label{tab:cohen_d}
\end{longtable}

\vfill
\section*{B.6 Notes on Statistical Power}

With $n=5$ seeds, the Wilcoxon signed-rank test has limited statistical power. The following observations address this limitation:

\begin{itemize}
    \item For the Easy-difficulty comparisons, all 5 seeds consistently favor the proposed method (W=0), giving strong non-parametric evidence despite the small $n$.
    \item The Hard-difficulty comparison (Proposed vs LOF) shows mixed results ($W=4$, $d=-0.68$), consistent with the acknowledged limitation that LOF outperforms on hard anomalies.
    \item Cohen's $d$ effect sizes provide a complementary measure that is meaningful even with small $n$; large effects (|d| > 0.8) indicate practical significance regardless of $p$-value.
    \item The block bootstrap CIs (Table~\ref{tab:bootstrap_ci}) are computed with 1000 replicates on the full dataset (3.08M records), providing strong statistical evidence independent of the seed-level comparisons.
    \item The McNemar test (Table~\ref{tab:mcnemar}) uses the full evaluation dataset ($n=3.08$M records), providing the strongest statistical evidence.
\end{itemize}
